# A full-parent terminal for a finite positive-d commuting edge

25 September 2026. New work; the round5 package is unchanged.

This theorem closes the nondegenerate finite-angle part of a distance-three
jet face when one of its two omitted events is allowed to open. It applies
to the C1 and C2 routing wherever the hypotheses below hold. It does not
assert completion of their projective endpoints or flat core boundaries.

## 1. Precise separate-field statement

Let the coefficients belong to the bounded valuation ring of a real closed
extension (K_c\supset\mathbb R), and let an independent measurement/state
test belong to the bounded valuation ring of (K_s\supset\mathbb R), with
projective charts used as necessary. Assume the coefficient vector is
physical over (K_c), is nonnegative, and satisfies **exactly**

\[
\lambda=0.
\]

Its real residue (c_0) is assumed to satisfy

\[
d_0>0,\qquad (e_2)_0=0.
\tag{1}
\]

There are no exact support requirements on (e_1,e_3,e_4,e_5,\alpha,\beta).
In particular any of them may open infinitesimally from zero. On the exact
residue face (\lambda=e_2=0), write the known determinant decomposition as

\[
Q_{c_0}(U,V)=U E(V)+U^2 B(V),
\quad E(V)=q_{10}+q_{11}V+q_{12}V^2,
\quad B(V)=q_{20}+q_{21}V+q_{22}V^2.
\tag{2}
\]

Fix a real (v_0\ne0), and suppose

\[
q_{10}>0,\quad q_{12}>0,\quad E(v_0^2)=0,\quad B(v_0^2)>0.
\tag{3}
\]

Also assume the three-state complement (M_0(v_0)) in (5) below is positive
definite. These are entirely real residue conditions. The latter condition
is automatic if (\alpha_0+(e_5)_0>0), and is displayed explicitly to avoid
an unproved support-exhaustion claim.

For every independent test with

\[
\bar u=0,\qquad \bar v=v_0,
\qquad [\bar W:\bar X:\bar Y:\bar Z]=[1:0:0:0],
\tag{4}
\]

the original full slack is a finite sum of squares in the original bounded
tensor ring (B_c\otimes_{\mathbb R}B_s). All coefficient-opening rates and
all test rates in this chart are allowed. No quotient by an infinitesimal
test quantity or mixed tensor element is used.

The same statement with (E(v_0^2)>0) is the immediate positive-definite
four-normal exit described in section 7.

## 2. Event normalization and the complete polynomial

Use the full parent in coefficient order

\[
(d,\lambda,e_1,e_2,e_3,e_4,e_5,\alpha,\beta).
\]

The event amplitude (R_2=vW-X+uvY-uZ) is a pure tested unit in (4).
Divide the state by (R_2). This is a legal pure-support unit operation and
sets

\[
R_2=1,\qquad W=(1+X-uvY+uZ)/v.
\]

The entire opened-event contribution is now the **constant** (4e_2).
Call the resulting full raw slack (F(X,Y,Z,u,v)). Its coefficients are
polynomials in the point coefficients and Laurent polynomials in (v).
The only denominators are powers of the pure tested unit (v).

At (u=0), its state Hessian divided by two is

\[
M_0(v)=
\begin{pmatrix}
2(d+2e_1+2\beta)(1+v^2)&0&0\\
0&2d(1+v^2)+4(e_3+e_4)(1+v^2)+4e_5+4\alpha v^2
 &4v(e_5-\alpha)\\
0&4v(e_5-\alpha)&4\alpha+4e_5v^2
\end{pmatrix}.
\tag{5}
\]

Its lower block is positive definite if (d>0) and

\[
\alpha d+2\alpha e_3+2\alpha e_4+2\alpha e_5(1+v^2)
+e_5v^2(d+2e_3+2e_4)>0.
\]

Introduce **coefficient-side center variables** (x,a,b,U) by the formal
substitution

\[
X=u^2x,\quad Y=ua,\quad Z=ub,\quad U=u^2.
\]

The exact full identity is

\[
F(u^2x,ua,ub,u,v)=4e_2+u^2K(x,a,b,v,U),
\quad U=u^2.
\tag{6}
\]

Here (K) is polynomial in (x,a,b,U), with Laurent-polynomial (v)
coefficients. More strongly,

\[
K_x=U E_x,\qquad
\partial_x E_x\big|_{U=0}=4(d+2e_1+2\beta)(1+v^2)>0,
\tag{7}
\]

and (K(x,a,b,v,0)) is independent of (x). These statements have been
verified directly on all 105 terms of (K), starting from all 106 terms of
the normalized full parent. The formal substitution defines the center;
it is never used to form (X/u^2,Y/u,Z/u) in the independent tensor.

## 3. The residue center and positive quartic coefficient

At the real residue, minimizing the positive state block in (a,b) gives
a scalar jet profile

\[
r(v)=\frac{16(1+v^2)^2}{v^2\det M_0(v)}E(v^2).
\tag{8}
\]

Physicality and (3) imply

\[
E(V)=q_{12}(V-v_0^2)^2.
\]

Consequently (r(v_0)=r'(v_0)=0), (r''(v_0)>0). The Hessian of

\[
K_0(a,b,v)=K(x,a,b,v,0)
\]

in ((a,b,v)) is positive definite at its unique local zero

\[
(a_0,b_0,v_0).
\]

This follows by completing the strictly positive ((a,b)) quadratic and
using (r''(v_0)>0), so it controls the full Hessian, not just its
determinant. Equation (E_x=0) then determines a unique (x_0) by (7).

Set (q=(x,a,b,v)). The Jacobian of

\[
\mathcal E(q,U)=(E_x,K_a,K_b,K_v)
\tag{9}
\]

in (q) at ((c_0,q_0,0)) is block upper triangular. Its top-left entry is
the positive number in (7), and its lower (3\times3) block is the positive
Hessian just proved. Thus it is invertible. The Nash implicit function
theorem gives a coefficient-only local branch (q_*(c,U)).

The optimized function

\[
k(c,U)=K(q_*(c,U),U)
\]

satisfies

\[
k(c_0,0)=0,\qquad
\gamma:=k_U(c_0,0)
=\frac{16(1+v_0^2)^2}{v_0^2\det M_0(v_0)}B(v_0^2)>0.
\tag{10}
\]

For clarity, (10) uses the whole parent determinant and the positive state
block. The event-normalized four-by-four matrix has determinant

\[
\det P/v^2.
\]

Its state Schur complement is its determinant divided by the positive
three-state determinant. At (e_2=0), (2) therefore gives (8), and at the
zero in (3) its next (U^2) coefficient is exactly (10). The equations
for (a,b) optimize the order-(U) state terms; the divided equation for

\(x\) optimizes the order-(U^2) state term. No inverse of this state block
is used in the final tensor identity; this real calculation only proves
strict positivity of the residue coefficient.

## 4. Two coefficient-only centers

Put (\kappa=k(c,0)). It is a pure coefficient infinitesimal, with either
sign allowed.

If (\kappa\ge0), take (q_c=q_*(c,0)) and (U_c=0). The ordinary
derivatives (K_a,K_b,K_v) vanish there, and (E_x=0).

If (\kappa<0), solve

\[
\mathcal E(q_c,U_c)=0,\qquad
K(q_c,U_c)+U_cK_U(q_c,U_c)=0.
\tag{11}
\]

The full Jacobian at the real base point is block upper triangular with
the invertible block (9) and final diagonal entry (2\gamma>0). The
derivatives of the final equation in (q) vanish at the base point.
Hence there is a unique nearby coefficient-only solution.

It satisfies (U_c>0). Indeed along the normal branch the last equation
is (g(U)=k(U)+Uk'(U)=0), with (g(0)=\kappa<0) and derivative having
positive residue (2\gamma). This argument is local and transfers to
the real closed coefficient field. It uses no comparison with the
independent tested (u).

Define

\[
\mu=4e_2+U_cK(q_c,U_c).
\tag{12}
\]

This is nonnegative: realize \(u_c=\sqrt{U_c}\), (v=v_c),

\[
(X,Y,Z)=(U_cx_c,u_ca_c,u_cb_c),
\]

and restore (W) by the event normalization. This is an actual
coefficient-side state and measurement pair, so full physicality applies.

## 5. Exact full-tail identity in the original independent tensor

Keep (q_c) fixed and put

\[
N_x=X-u^2x_c,\quad N_y=Y-ua_c,\quad N_z=Z-ub_c,
\quad T=v-v_c,\quad N_v=uT.
\tag{13}
\]

All these belong to the original bounded tensor ring. No small quantity
has been divided out.

The state dependence of the normalized full parent is exactly quadratic.
Consequently its entire expansion is

\[
F=4e_2+U K(q_c\text{ with }v\text{ variable},U)
+U E_xN_x+uK_aN_y+uK_bN_z+N_s^TMN_s,
\quad N_s=(N_x,N_y,N_z),
\tag{14}
\]

where the derivative functions on the right are evaluated at

\((x_c,a_c,b_c,v,U)\). The independent checker verifies this identity,
all three linear coefficients, and the complete quadratic remainder.

At (v=v_c,U=U_c) all four center equations vanish. Exact difference
division therefore writes the linear terms as combinations of

\[
N_x U(U-U_c),\quad N_yu(U-U_c),\quad N_zu(U-U_c),
\quad N_xuN_v,\quad N_yN_v,\quad N_zN_v.
\tag{15}
\]

The coefficient functions are bounded. Laurent denominators involve only
the pure units (v,v_c); differences of their powers are divided using
finite algebraic identities. Taylor division of the profile in (v)
gives an additional term (N_v^2h) and a cross term

\[
N_vu(U-U_c)h_1.
\tag{16}
\]

Thus all cross terms with the last scalar generator have an extra (u)
or (U), and have zero residue at every independent test rate.

For (\kappa\ge0), univariate division at (U=0) gives the exact form

\[
F=4e_2+\kappa u^2+G^T H G,
\quad G=(N_x,N_y,N_z,N_v,U).
\tag{17}
\]

For (\kappa<0), stationarity in (11) and (12) gives a double zero of
the fixed-center univariate profile at (U_c). Monic polynomial division
(its coefficients may contain only the pure coefficient unit (v_c))
therefore gives

\[
F=\mu+G^THG,
\quad G=(N_x,N_y,N_z,N_v,U-U_c).
\tag{18}
\]

In both cases the residue Gram is positive definite. Its first block is
the positive state-(X) entry together with one half of the positive
Hessian of (K_0(a,b,v)); its last entry is (\gamma>0). The cross entries
to the last generator vanish by (15)–(16). This statement also accounts
for the possibly nonzero residue cross terms (N_yN_v,N_zN_v), which are
part of the positive full jet Hessian.

The bounded positive-residue Gram lemma from round5 now makes

\(G^THG\) an SOS in (B_c\otimes B_s). The remaining coefficients
\(e_2,\kappa,\mu\) are nonnegative pure coefficient quantities in their
respective branches and hence squares. Finally multiply the roots by the
original pure tested event amplitude to undo normalization. If normalized
physical slack rather than raw slack is desired, the norm and measurement
denominators are positive pure tested units on this chart and can likewise
be restored. This proves section 1.

## 6. Nonvacuity and exact checks

The real residue

\[
(d,\lambda,e_1,e_2,e_3,e_4,e_5,\alpha,\beta)
=(2,0,0,0,1,1,1,1,0)
\]

is physical, lies in both the C1 and C2 omission masters, and has

\[
Q=16U\{4UV+3(V-1)^2\}.
\]

At (v_0=1), the center is

\[
(x_0,a_0,b_0)=(-1/2,1/2,1),\qquad \gamma=2.
\]

The divided (x) pivot is 16, and the full lower Hessian is

\[
\begin{pmatrix}64&0&48\\0&16&16\\48&16&64\end{pmatrix}\succ0.
\]

`verify_positive_d_finite_edge.py` rederives the full parent, event
normalization, divided center equations, exact linear jets, and every
state-quadratic tail term. Its report is
`POSITIVE_D_FINITE_EDGE_CHECKS.json`: **15 identities, 30 scalar entries,
all passed**. It includes the displayed example and both independent
calculations of its positive quartic coefficient.

## 7. Exact scope of the remaining cases

If (E(v_0^2)>0), the original four-normal expansion in

\((X,Y,Z,u)\) after event normalization has a positive residue Gram.
The term (4e_2) remains separate, so this strict finite-angle exit needs
no critical center. If the tested state has positive residue loss, the
scalar bounded-positive-residue lemma applies immediately.

The theorem above treats a finite nonzero zero of a nonzero quadratic
edge polynomial with both endpoint coefficients positive. It leaves:

- \(v_0=0\) or (v_0=\infty), where the normalization or the jet Hessian
  changes;
- flat outside/core cases (E\equiv0), where the (v) Hessian vanishes;
- distance-four corner jets that are not relabelings of this finite-edge
  configuration;
- coefficient charts in which both events of the chosen distance-three
  omission pair are nonzero. For applications to C1/C2 extreme routing,
  a permanently omitted event must first be included in the chosen pair.

No density, continuity of certificates, or unproved global gluing is used
to include those cases.
